% ============================================================
% Phase 6: Circular Dependency Resolution
% (Stage 1 ↔ Stage 2 as Filippov Hybrid System)
% Status: FULL RECONSTRUCTION
% ============================================================

% --- Definitions and Assumptions ---
\begin{definition}[Circular Dependency]
Stage~1 (depth estimation) requires $\tilde{\sigma}_A$ as input to
compute the fusion signal $g_A$; Stage~2 (schema coherence estimation)
requires $\delta_A$ as input to compute the ground-truth proxy. This is
a mutual dependency:
\[
\delta_A = F(\tilde{\sigma}_A), \qquad \tilde{\sigma}_A = G(\delta_A),
\]
where $F: \mathbb{R} \to \mathbb{R}$ and $G: \mathbb{R} \to \mathbb{R}$
are the Stage~1 depth estimator and Stage~2 schema estimator respectively,
each defined on the compact domain
$\mathcal{D} = [0, \delta_{\max}] \times [0, 1] \subset \mathbb{R}^2$.
\end{definition}

\begin{definition}[Piecewise-Smooth Hybrid System]
A Filippov system is a differential equation with a piecewise-smooth
right-hand side:
\[
\dot{x} = f_i(x), \quad x \in S_i,
\]
where $\{S_i\}$ is a finite partition of the state space and $f_i$ is
smooth on each $S_i$. On the switching boundary $\Sigma_{ij} = \bar{S}_i
\cap \bar{S}_j$, the vector field is defined by Filippov convex
combination.
\end{definition}

\begin{assumption}[Lipschitz Continuity]
The Stage~1 estimator $F$ is $L_F$-Lipschitz and the Stage~2 estimator
$G$ is $L_G$-Lipschitz on the domain $\mathcal{D}$:
\begin{align*}
|F(\tilde{\sigma}_1) - F(\tilde{\sigma}_2)| &\leq L_F |\tilde{\sigma}_1 - \tilde{\sigma}_2|, \\
|G(\delta_1) - G(\delta_2)| &\leq L_G |\delta_1 - \delta_2|,
\end{align*}
with Lipschitz constants $L_F, L_G \geq 0$.
\end{assumption}

\begin{assumption}[Contraction Condition]
The product of Lipschitz constants satisfies $L_F L_G < 1$, making
the composition $G \circ F$ a contraction on $\mathcal{D}$.
\end{assumption}

% --- Proposition Statement ---
\begin{proposition}[Circular Dependency Resolution]
\label{prop:hybrid}
Under Assumptions~1--2, the Stage~1/2 circular dependency admits a
unique consistent solution via the Filippov hybrid system formulation,
with the mutual estimation converging after at most
$K = \lceil \log(\theta_{\text{conv}} / d_0) / \log(L_F L_G) \rceil$
alternating rounds, where $d_0$ is the initial distance and
$\theta_{\text{conv}}$ the convergence tolerance.
\end{proposition}

% --- Proof ---
\begin{proof}

\paragraph{Step 1: Lipschitz bounds for $F$ and $G$.}
The Stage~1 depth estimator $F$ computes $\delta_A$ from the coupled
ODE system (Eqs.~14, 28, 29) with $\tilde{\sigma}_A$ as the schema
input. The ODE solution is differentiable with respect to the initial
condition and parameters (Proposition~3.1). By standard sensitivity
analysis for ODEs (the variation-of-constants formula), the derivative
$\partial \delta_A / \partial \tilde{\sigma}_A$ is bounded on the
compact interval $[0, T]$ by:
\[
\left|\frac{\partial \delta_A}{\partial \tilde{\sigma}_A}\right|
\leq \frac{\lambda_c \delta_{\max}}{\min(\lambda_c, \rho, \gamma)}
\cdot e^{(\eta_{\max} + \lambda_c) T} \triangleq L_F,
\]
where $\delta_{\max}$ is the maximum depth (Proposition~3.2), and the
factor $\lambda_c \delta_{\max}$ bounds the coupling term
$\partial \dot{\delta}_A / \partial \tilde{\sigma}_A = \lambda_c \gamma_\sigma \delta_A (1 - r_A)$.

The Stage~2 schema estimator $G$ computes $\tilde{\sigma}_A$ as a
weighted fusion of three proxy signals (GCA, RGA, AC; Section~3.1.4).
Each signal is bounded between 0 and 1, and the fusion is a convex
combination with weights $w_g + w_r + w_c = 1$. The mapping from
$\delta_A$ to $\tilde{\sigma}_A$ is therefore 1-Lipschitz when the
proxy signals themselves are 1-Lipschitz in $\delta_A$:
\[
\left|\frac{\partial \tilde{\sigma}_A}{\partial \delta_A}\right|
\leq w_g \cdot L_{\text{GCA}} + w_r \cdot L_{\text{RGA}} + w_c \cdot L_{\text{AC}}
\leq L_G,
\]
where $L_{\text{GCA}}$, $L_{\text{RGA}}$, $L_{\text{AC}}$ are the
Lipschitz constants of the individual proxy signal extractors.
In practice, each proxy is computed as a bounded differentiable
function of the hidden states, ensuring $L_G < \infty$.

\paragraph{Step 2: Contraction of $G \circ F$.}
By the chain rule and Assumption~1:
\[
\left|\frac{d(G \circ F)}{d\tilde{\sigma}_A}\right|
= \left|\frac{dG}{d\delta_A} \cdot \frac{dF}{d\tilde{\sigma}_A}\right|
\leq L_G \cdot L_F.
\]
Under Assumption~2 ($L_F L_G < 1$), the Banach fixed-point theorem
guarantees that the iteration:
\[
\tilde{\sigma}_A^{(k+1)} = G(F(\tilde{\sigma}_A^{(k)}))
\]
converges to a unique fixed point $\tilde{\sigma}_A^*$ from any
initial $\tilde{\sigma}_A^{(0)} \in [0, 1]$.

\paragraph{Step 3: Convergence rate.}
The contraction mapping theorem provides a geometric convergence rate:
\[
\|\tilde{\sigma}_A^{(k)} - \tilde{\sigma}_A^*\|
\leq (L_F L_G)^k \cdot \|\tilde{\sigma}_A^{(0)} - \tilde{\sigma}_A^*\|.
\]
For a tolerance $\theta_{\text{conv}} > 0$, the number of iterations
required is:
\[
K = \lceil \log(\theta_{\text{conv}} / d_0) / \log(L_F L_G) \rceil,
\]
where $d_0 = \|\tilde{\sigma}_A^{(0)} - \tilde{\sigma}_A^*\| \leq 1$
is the initial distance.

For typical parameter values:
\begin{itemize}[nosep]
\item $L_F \approx 0.5$ (moderate ODE sensitivity, bounded training horizon)
\item $L_G \approx 0.6$ (convex fusion of 1-Lipschitz proxy signals)
\item $L_F L_G \approx 0.3$, giving $(0.3)^k$ convergence
\item For $\theta_{\text{conv}} = 0.01$ and $d_0 = 1$: $K = \lceil \log(0.01) / \log(0.3) \rceil = \lceil 3.82 \rceil = 4$ rounds
\end{itemize}
The 2--3 round convergence observed in numerical experiments
(Appendix~C) corresponds to the tighter bound achieved when
$L_F L_G \approx 0.1$ (achieved when $T$ is small, limiting the
exponential factor in $L_F$).

\paragraph{Step 4: Hybrid system formulation.}
The three discrete modes resolve the circular dependency without
simultaneous equation solving:

\begin{enumerate}[nosep]
\item \textbf{Mode S1 (Stage~1 active):} Use prior $\tilde{\sigma}_A^{(k)}$
  to solve the depth ODE and produce $\delta_A^{(k)} = F(\tilde{\sigma}_A^{(k)})$.
\item \textbf{Mode S2 (Stage~2 active):} Use $\delta_A^{(k)}$ to compute
  $\tilde{\sigma}_A^{(k+1)} = G(\delta_A^{(k)})$.
\item \textbf{Mode F (Fixed-point):} If
  $\|\tilde{\sigma}_A^{(k+1)} - \tilde{\sigma}_A^{(k)}\| < \theta_{\text{conv}}$,
  output $(\delta_A^*, \tilde{\sigma}_A^*)$.
\end{enumerate}

The mode transition function $\phi: \mathbb{R}^2 \to \{\text{S1}, \text{S2}, \text{F}\}$ is:
\[
\phi(\tilde{\sigma}_A, \delta_A) = \begin{cases}
\text{S1} & \text{if } \|\tilde{\sigma}_A^{(k+1)} - \tilde{\sigma}_A^{(k)}\| > \theta_{\text{conv}}
\text{ and current mode is S2}, \\
\text{S2} & \text{if } \|\delta_A^{(k+1)} - \delta_A^{(k)}\| > \theta_{\text{conv}}
\text{ and current mode is S1}, \\
\text{F} & \text{otherwise}.
\end{cases}
\]

This is a Filippov system with two switching manifolds:
\begin{align*}
\Sigma_{12} &= \{(\tilde{\sigma}_A, \delta_A) : \|\tilde{\sigma}_A - G(\delta_A)\| = \theta_{\text{conv}}\}, \\
\Sigma_{21} &= \{(\tilde{\sigma}_A, \delta_A) : \|\delta_A - F(\tilde{\sigma}_A)\| = \theta_{\text{conv}}\},
\end{align*}
where on each manifold the vector field is the convex combination of
the adjacent modes' fields, ensuring existence of solutions by
Filippov's theorem.

\paragraph{Step 5: Uniqueness of the fixed point.}
By the Banach fixed-point theorem, $G \circ F$ has exactly one fixed
point $\tilde{\sigma}_A^*$ on the compact domain $[0, 1]$. The
corresponding depth is $\delta_A^* = F(\tilde{\sigma}_A^*)$, which is
unique because $F$ is Lipschitz (hence single-valued). The pair
$(\tilde{\sigma}_A^*, \delta_A^*)$ is the unique consistent solution
to the circular dependency. $\square$
\end{proof}

% --- Boundary Cases ---
\begin{boundary-case}
\textbf{Cold start ($\tilde{\sigma}_A^{(0)} = 0$):}
The initial distance $d_0 = \|\tilde{\sigma}_A^*\| \leq 1$ is maximal.
The required iteration count $K = \lceil \log(\theta_{\text{conv}}) / \log(L_F L_G) \rceil$
increases by approximately $\lceil \log(1/d_0) / \log(L_F L_G) \rceil$ rounds.
For $L_F L_G = 0.3$ and $\theta_{\text{conv}} = 0.01$, cold start
requires $K = 4$ rounds vs.\ $K = 2$ for warm start with
$d_0 = 0.01$.

\textbf{$L_F L_G \geq 1$ (not a contraction):}
The fixed-point iteration may diverge or oscillate. In this regime,
the hybrid formulation is replaced by a simultaneous equation solver
(Newton--Raphson on the coupled system $(\tilde{\sigma}_A, \delta_A)$).
The Newton--Raphson iteration converges quadratically provided the
Jacobian $\mathbf{I} - \mathbf{J}_{G \circ F}$ is non-singular at the
solution.

\textbf{Time-discretized implementation:}
The practical pipeline (Algorithm~3.2) uses alternating updates each
epoch rather than inner-loop iteration. This single-step alternating
scheme is equivalent to one iteration of the fixed-point map per
epoch. The convergence analysis holds asymptotically as the epoch
count grows, because the composition $G \circ F$ is applied
repeatedly across epochs. The effective contraction rate per epoch is
$(L_F L_G)^{1/\tau_{\text{epoch}}}$ where $\tau_{\text{epoch}}$ is
the number of ODE integration steps per epoch, making the practical
convergence slower but still geometric.

\textbf{Discontinuous proxy signals:}
If any proxy signal has discontinuities (e.g., threshold-based metrics),
the Lipschitz assumption (Assumption~1) is violated. The Filippov
formulation still guarantees existence of solutions via the convex
combination on switching boundaries, but the convergence rate bound
no longer holds. In practice, all three proxies (GCA, RGA, AC) are
continuous functions of the hidden state activations, so
discontinuity is not expected.
\end{boundary-case}

% --- Circular Dependency Check ---
\begin{circularity-check}
The proof uses the following external results:
\begin{itemize}[nosep]
\item Proposition~3.1 (local existence and uniqueness of ODE solutions)
  to justify differentiability of $F$ with respect to parameters.
\item Proposition~3.2 (boundedness) to define $\delta_{\max}$ and the
  compact domain $\mathcal{D}$.
\item The ODE definitions from Section~3.2 and the proxy signal
  definitions from Section~3.1.4.
\item Standard results from functional analysis: Banach fixed-point
  theorem, Filippov's existence theorem for discontinuous differential
  equations.
\end{itemize}

\textbf{Self-circularity:} The Lipschitz constants $L_F$ and $L_G$
are stated as bounded but not explicitly computed in closed form.
The bound on $L_F$ depends on the ODE integration horizon $T$ and
the parameter bounds, which are themselves derived from
Proposition~3.2 (boundedness). This is not circular: Proposition~3.2
establishes the invariant set independently of the Stage~1/2
estimation procedure.

\textbf{No dependence on later results:} The proof does not use any
proposition or theorem from Section~4 (Cognitive Dimensions) or
later. It depends only on Section~3 infrastructure (ODE definitions,
existence, boundedness).

\textbf{No dependence on Proposition~3.4 (equilibria):} The equilibrium
analysis is not required because $F$ and $G$ are defined as
finite-horizon ODE solutions and proxy signal evaluations,
respectively, neither of which requires reaching equilibrium.
\end{circularity-check}
